\section{问题一：预热阶段的径向热湿传递}
\subsection{控制方程的建立}
\medskip
\noindent\textbf{（1）基于热量守恒的径向导热模型}

烘干过程中，烘房通过药材表面向内部传递热量，使药材各位置的温度随时间变化。根据轴对称假设，忽略轴向和周向的温度差异，将温度表示为径向位置与时间的函数$T(r,t)$。在药材内部取半径为$r$至$r+\dd r$、长度为$L$的薄圆柱环层，其体积为
\[
\dd V=\pi\bigl[(r+\dd r)^2-r^2\bigr]L
\approx2\pi rL\,\dd r.
\]

以径向向外为热流的正方向。根据傅里叶导热定律，单位面积上的热流密度与温度梯度成正比、方向相反，即
\[
q(r,t)=-k\frac{\partial T}{\partial r},
\]
其中，$k$为药材的导热系数。半径为$r$的圆柱面面积为$2\pi rL$，因此通过该面的热流率为
\[
Q(r,t)=2\pi rLq(r,t)
=-2\pi rLk\frac{\partial T}{\partial r}.
\]

在不计内部热源、汽化潜热及水分迁移携带热量的条件下，环层内的蓄热速率等于内外两侧的净流入热流率。由热量平衡可得
\[
\rho c_p\,\dd V\frac{\partial T}{\partial t}
=Q(r,t)-Q(r+\dd r,t).
\]
将右端按$r$作一阶展开，并代入环层体积，得到
\[
\rho c_p(2\pi rL\,\dd r)\frac{\partial T}{\partial t}
=-\frac{\partial}{\partial r}(2\pi rLq)\,\dd r.
\]
进一步代入傅里叶定律，约去公共因子，得到药材内部的径向热传导方程
\begin{equation}\label{eq_heat1}
\rho c_p\frac{\partial T}{\partial t}
=\frac1r\frac{\partial}{\partial r}\left(rk\frac{\partial T}{\partial r}\right),
\qquad 0<r<R.
\end{equation}
问题一采用常数热物性参数$\rho=820$ kg/m$^3$、$c_p=2600$ J/(kg$\cdot$K)、$k=0.36$ W/(m$\cdot$K)。令热扩散率$\alpha=k/(\rho c_p)$，则式\eqref{eq_heat1}可化为
\[
\frac{\partial T}{\partial t}
=\alpha\left(\frac{\partial^2T}{\partial r^2}
+\frac1r\frac{\partial T}{\partial r}\right).
\]
其中，含$1/r$的项反映了圆柱面传热面积随半径变化的几何特征。烘房对药材的加热作用通过外表面换热边界引入，内部温度分布则由上述方程确定。

\medskip
\noindent\textbf{（2）基于水分守恒的非线性扩散模型}

药材表面与烘房发生水分交换后，内部逐渐形成含水率梯度，水分随之由内部向表面迁移。题目中的$C(r,t)$为干基含水率，即单位干物质质量所含的水质量。为建立水质量守恒关系，按照假设（5），将单位体积中的干物质质量记为常数$s_0$，则单位体积中的水质量为$s_0C$。

采用有效Fick扩散关系描述内部水分迁移。以径向向外为正方向，相对干物质骨架的水质量通量为
\[
j(r,t)=-s_0D(C)\frac{\partial C}{\partial r},
\]
其中，$D(C)$为有效水分扩散系数，单位为m$^2$/s。当含水率由内向外降低时，$\partial C/\partial r<0$，对应$j>0$，即水分向表面迁移。

沿用热传导模型的环层收支思路，水质量的变化率等于净流入的水质量流率。由于环层体积与参考干物质密度保持不变，有
\[
s_0(2\pi rL\,\dd r)\frac{\partial C}{\partial t}
=-\frac{\partial}{\partial r}(2\pi rLj)\,\dd r.
\]
代入扩散通量表达式，消去$s_0$和公共几何因子，得到
\begin{equation}\label{eq_moist1}
\frac{\partial C}{\partial t}
=\frac1r\frac{\partial}{\partial r}\!\left[rD(C)\frac{\partial C}{\partial r}\right],\qquad
D(C)=7\times10^{-9}\exp\!\left(-\frac{0.89}{C}\right).
\end{equation}
与常系数扩散不同，式\eqref{eq_moist1}中的$D$随局部含水率变化，因而水分迁移过程具有非线性。将其展开可得
\[
\frac{\partial C}{\partial t}
=D(C)\left(\frac{\partial^2C}{\partial r^2}+\frac1r\frac{\partial C}{\partial r}\right)
+D'(C)\left(\frac{\partial C}{\partial r}\right)^2,
\qquad D'(C)=\frac{0.89D(C)}{C^2}>0.
\]
因此，扩散系数不能作为常数直接移出微分算子。后续求解保留式\eqref{eq_moist1}的通量形式，以便在离散过程中统一计算相邻环层之间的水分交换，并检验整体水分收支。

综上，问题一的热物性均为常数，水分扩散系数仅依赖含水率。在不显式考虑汽化潜热及水分携热的假设下，温度方程与水分方程相互独立，可分别结合各自的初始条件和表面交换条件进行求解。

\subsection{初始条件与边界条件}
入炉时药材各处温度和含水率均匀，题给初始状态为
\begin{equation}\label{eq_initial}
T(r,0)=28,\qquad C(r,0)=2.55,\qquad 0\le r\le R.
\end{equation}
式\eqref{eq_initial}分别以$^\circ$C和kg/kg为单位。根据径向对称性，中心处满足
\[
\frac{\partial T}{\partial r}(0,t)
=\frac{\partial C}{\partial r}(0,t)=0.
\]
这一条件表示中心温度和含水率的径向梯度为零；中心状态仍随内部传递过程变化。

在外表面，导热通量与对流换热通量相平衡，水分扩散通量与假设（4）中的等效水分交换通量相平衡。因此采用Robin边界
\begin{equation}\label{eq_robin}
\begin{aligned}
-kT_r(R,t)&=h\,[T_s-T_a(t)],\\
-D(C_s)C_r(R,t)&=h_m\,[C_s-C_a(t)].
\end{aligned}
\end{equation}
其中，$T_s=T(R,t)$、$C_s=C(R,t)$，$h=25$ W/(m$^2\cdot$K)、$h_m=8\times10^{-7}$ m/s。式\eqref{eq_robin}中，$T_a>T_s$时向外热流为负，表示热量进入药材；$C_s>C_a$时水分向外迁移。表面温度和含水率由内部传递与外部交换共同决定，因而通常与烘房状态不同。

环境函数$T_a(t)$和$C_a(t)$沿用“输入整理与时间延续”中的PCHIP处理，按式\eqref{eq_pchip}对附件1前1800 s的31组观测插值，在求解器所需时刻计算边界输入。这里的$h_m(C_s-C_a)$仍按等效交换关系使用，$C_a$的含义与模型假设保持一致。

均匀初始水分使表面初始梯度为零，而交换边界在接触烘房后立即产生失水通量，二者在$t=0$角点不相容。本文保留题给初值，在$t>0$施加表面交换，求得初始表面层的瞬态发展，并把前1 s的表面结果纳入网格检验。

\subsection{模型的数值求解}
\medskip
\noindent\textbf{（1）空间离散}

预热初期，药材表面含水率变化较快，较陡的径向梯度主要出现在外层。为分辨这一变化，采用节点型有限体积法，在靠近表面处逐步加密网格，并使题目规定的0、0.1、$\ldots$、2.0 cm输出位置均为实际节点。将每个0.1 cm基本区间细分，在0至1.5、1.5至1.8、1.8至1.9、1.9至2.0 cm四段内，分别细分为$8m$、$16m$、$32m$、$64m$份，记所得网格为$G_m$，共有$264m+1$个节点。本问采用$G_8$，即2113个节点，并通过加密比较检验空间离散误差。

设全部节点依次为$r_i$（$i=0,\ldots,N$），内部面为$r_{i+1/2}=(r_i+r_{i+1})/2$。控制体左右端点为$a_i,b_i^*$，中心端取0、表面端取$R$，定义
\begin{equation}\label{eq_weights}
W_i=\int_{a_i}^{b_i^*}r\dd r=\frac{(b_i^*)^2-a_i^2}{2},
\qquad \sum_{i=0}^NW_i=\frac{R^2}{2}.
\end{equation}
真实控制体体积为$2\pi LW_i$。式\eqref{eq_weights}保留了不同半径处环层体积的差异，中心与表面处的控制体均在物理边界截取。

\medskip
\noindent\textbf{（2）通量处理与半离散方程}

记$D_i=D(C_i)$，在相邻节点之间对扩散系数取调和平均，得到已乘半径的向外通量
\begin{equation}\label{eq_flux}
\begin{aligned}
F^T_{i+1/2}&=-r_{i+1/2}k\frac{T_{i+1}-T_i}{r_{i+1}-r_i},\\
F^C_{i+1/2}&=-r_{i+1/2}\frac{2D_iD_{i+1}}{D_i+D_{i+1}}
\frac{C_{i+1}-C_i}{r_{i+1}-r_i}.
\end{aligned}
\end{equation}
中心控制体在$r=0$处的边界通量取零，表面通量直接取$Rh(T_N-T_a)$、$Rh_m(C_N-C_a)$。对式\eqref{eq_heat1}和式\eqref{eq_moist1}作控制体积分，并用节点状态近似体平均，得到
\begin{equation}\label{eq_fvm}
\rho c_pW_i\dot T_i=F^T_{i-1/2}-F^T_{i+1/2},\qquad
W_i\dot C_i=F^C_{i-1/2}-F^C_{i+1/2}.
\end{equation}
式\eqref{eq_flux}在相邻控制体中使用同一内部面通量，使其在式\eqref{eq_fvm}中以相反符号出现，对所有控制体求和后仅保留表面交换项。

中心处采用控制体积分，避免直接计算$1/r$。若首个节点间距为$\delta=r_1-r_0$，则$W_0=\delta^2/8$，相应温度变化率为$\dot T_0=4\alpha(T_1-T_0)/\delta^2$，水分方程具有同样的几何系数。表面节点则保留自身蓄积项，同时接收内部面通量和外部交换通量，从而直接求得真实表面状态。

\medskip
\noindent\textbf{（3）时间推进与结果输出}

空间离散后，两场分别形成常微分方程组，统一记为$\dot{\boldsymbol y}=\boldsymbol f(t,\boldsymbol y)$，其中$\boldsymbol y$分别取全部节点的温度或含水率。表面加密使显式方法受到较小稳定步长的限制，因此采用隐式后向差分公式（BDF）进行时间积分~\cite{bdf}。

以等步长二阶BDF为例，下一时刻的状态应满足
\begin{equation}\label{eq_bdf_residual}
\boldsymbol{\mathcal R}(\boldsymbol y^{n+1})
=3\boldsymbol y^{n+1}-4\boldsymbol y^n+\boldsymbol y^{n-1}
-2\Delta t\,\boldsymbol f(t_{n+1},\boldsymbol y^{n+1})
=\boldsymbol0.
\end{equation}
式\eqref{eq_bdf_residual}中的历史状态$\boldsymbol y^n$、$\boldsymbol y^{n-1}$已知，但通量也依赖待求的$\boldsymbol y^{n+1}$，因而需要求解隐式代数方程组。对水分场，$D(C)$随未知含水率变化，方程组具有非线性。计算时先由历史状态预测新状态，再进行牛顿型迭代：按试探状态更新扩散系数和通量，计算方程残差，求解线性修正方程并更新状态，直至满足收敛要求。

步内迭代收敛后，还需估计局部时间误差；误差满足容差时才接受该时间步，否则缩小步长重算。实际调用SciPy的变步长、变阶BDF实现，式\eqref{eq_bdf_residual}仅用于说明其隐式离散形式，计算并不固定采用二阶公式。

本问相对容差设为$10^{-10}$，温度与含水率绝对容差分别为$10^{-10}$、$10^{-12}$，最大步长为0.5 s。温度场与水分场分别积分；水分方程的每次右端函数求值均按当前含水率更新$D(C)$，线性修正方程的求解利用相邻节点耦合形成的稀疏结构。

在0至1800 s内完成积分后，利用求解器的连续输出提取题定时刻的状态，并读取相应径向节点。完整结果按1 s、0.1 cm间隔保存，题定表格从中选取指定时刻与位置。每秒输出是数据保存要求，求解器的内部步长仍由误差控制自适应选取。

\subsection{计算结果与分析}
表\ref{tab_q1T}、表\ref{tab_q1C}给出题定时刻与位置的温度和含水率，图\ref{fig_q1}展示相应径向剖面，用于比较预热过程中不同位置的升温与失水情况。
\begin{table}[htbp]\centering\small
\caption{30分钟内药材的温度（$^\circ$C）}\label{tab_q1T}
\renewcommand{\arraystretch}{1.12}
\begin{tabular}{cccccc}\toprule
 & \multicolumn{5}{c}{到药材中心的距离/cm}\\
时间/s & 0 & 0.5 & 1.0 & 1.5 & 2.0\\\midrule
100 & 28.0000 & 28.0003 & 28.0038 & 28.0319 & 28.1795\\
300 & 28.0406 & 28.0634 & 28.1513 & 28.3673 & 28.8486\\
600 & 28.4533 & 28.5361 & 28.8039 & 29.3155 & 30.1641\\
900 & 29.3243 & 29.4583 & 29.8757 & 30.6164 & 31.7295\\
1200 & 30.5429 & 30.7099 & 31.2225 & 32.1130 & 33.4295\\
1500 & 31.9960 & 32.1871 & 32.7667 & 33.7476 & 35.1211\\
1800 & 33.5758 & 33.7724 & 34.3642 & 35.3626 & 36.7879\\
\bottomrule\end{tabular}
\end{table}

\begin{table}[htbp]\centering\small
\caption{30分钟内药材的水分浓度（kg/kg）}\label{tab_q1C}
\renewcommand{\arraystretch}{1.12}
\begin{tabular}{cccccc}\toprule
 & \multicolumn{5}{c}{到药材中心的距离/cm}\\
时间/s & 0 & 0.5 & 1.0 & 1.5 & 2.0\\\midrule
100 & 2.5500 & 2.5500 & 2.5500 & 2.5500 & 2.2470\\
300 & 2.5500 & 2.5500 & 2.5500 & 2.5492 & 2.0508\\
600 & 2.5500 & 2.5500 & 2.5500 & 2.5353 & 1.8770\\
900 & 2.5500 & 2.5500 & 2.5497 & 2.5045 & 1.7547\\
1200 & 2.5500 & 2.5500 & 2.5482 & 2.4646 & 1.6586\\
1500 & 2.5500 & 2.5499 & 2.5445 & 2.4206 & 1.5787\\
1800 & 2.5500 & 2.5497 & 2.5383 & 2.3755 & 1.5102\\
\bottomrule\end{tabular}
\end{table}

\figpair{q1_T_profiles.pdf}{q1_C_profiles.pdf}{预热阶段温度与含水率径向剖面。外层先升温、先失水，中心含水率在30 min内变化很小；不同曲线对应题定输出时刻。}{fig_q1}

\FloatBarrier

温度分布体现了由外向内的升温过程。至1800 s，表面温度为36.7879 $^\circ$C，中心温度为33.5758 $^\circ$C，两者相差3.2121 $^\circ$C；表面仍低于同期烘房温度41.5130 $^\circ$C。热量先经表面对流进入药材，再通过内部导热向中心传播，有限的表面换热和内部导热共同造成了这一温度滞后。

水分变化主要集中在外层。末时刻表面含水率降至1.5102 kg/kg，而中心未舍入值为2.5499924 kg/kg，保留四位小数仍为2.5500 kg/kg。表面持续向烘房排湿，首先形成含水率梯度；在前30 min内，内部水分向外迁移尚不充分，因此中心变化较小，径向剖面表现为中心高、表面低。

含水率下降还会改变后续迁移速率。由$D'(C)>0$可知，失水使局部扩散系数减小。本问$D$由初始$4.9377\times10^{-9}$降至局部最小$3.8830\times10^{-9}$ m$^2$/s，表明在相同含水率梯度下，较干区域的水分传递能力有所降低。与冻结初始$D$的计算相比，含水率最大差达到0.0419 kg/kg，因而结果分析和正式求解均需保留$D(C)$的变化。

\subsection{模型检验}
\label{sec_q1_validation}
本节从空间与时间收敛、独立参考解和收支一致性三个方面检验第一问的计算可靠性。需要说明的是，这些检验验证的是模型方程、有限体积离散和BDF求解实现；径向近似、忽略蒸发潜热及等效水分边界等物理假设不能由本节单独证明，其适用性将在后文统一讨论。

\medskip
\noindent\textbf{（1）空间与时间收敛性检验}

依次采用$G_2$、$G_4$、$G_8$三级网格，节点数分别为529、1057和2113。在相同的1800个输出时刻及21个径向位置比较相邻两级解，温度与含水率的最大差及其发生位置见表\ref{tab_q1_convergence}。由$G_2/G_4$加密至$G_4/G_8$后，温度最大差由$2.289\times10^{-6}$ $^\circ$C降至$5.895\times10^{-7}$ $^\circ$C，含水率最大差由$1.828\times10^{-5}$ kg/kg降至$4.564\times10^{-6}$ kg/kg，两者均呈现随网格加密而减小的趋势。

\begin{table}[htbp]\centering\small
\caption{问题一空间与时间加密结果}\label{tab_q1_convergence}
\begin{tabular}{lccp{6.2cm}}\toprule
比较项 & $\max|\Delta T|$ & $\max|\Delta C|$ & 最大差位置\\[-2pt]
 & ($^\circ$C) & (kg/kg) & \\\midrule
$G_2/G_4$ & $2.289\times10^{-6}$ & $1.828\times10^{-5}$ & $T$：684 s，1.5 cm；$C$：1 s，2.0 cm\\
$G_4/G_8$ & $5.895\times10^{-7}$ & $4.564\times10^{-6}$ & $T$：678 s，1.5 cm；$C$：1 s，2.0 cm\\
$G_8$时间收紧 & $2.874\times10^{-8}$ & $1.179\times10^{-9}$ & $T$：1501 s，2.0 cm；$C$：1 s，2.0 cm\\\bottomrule
\end{tabular}
\end{table}

含水率最大差出现在入炉后的第1秒且位于表面，说明检验覆盖了初边值不相容所造成的最陡表面层。进一步在生产网格上收紧时间容差和最大步长后，温度与含水率最大差分别为$2.874\times10^{-8}$ $^\circ$C和$1.179\times10^{-9}$ kg/kg，均小于对应的空间加密差，表明时间离散误差不是当前主要误差来源。上述差值反映相邻计算配置之间的稳定程度，不直接等同于连续方程的真解误差。

题定7组时刻、5个位置的温度和含水率在$G_4/G_8$及时间收紧比较下均保留四位一致。完整逐秒、21个径向位置的输出中仍有少量数值处于四位舍入边界，因此本文将“四位稳定”限定为题定表格，不将其扩展为全部输出均已冻结的严格舍入证明。

\medskip
\noindent\textbf{（2）Bessel--Duhamel温度参考解}

问题一的热物性为常数，且在本问不显式考虑潜热的基线下，温度方程与水分方程相互独立，因此可以构造独立于径向有限体积离散的温度参考解。记$J_0,J_1$为第一类Bessel函数，对圆柱径向导热方程作分离变量，模态特征值满足
\begin{equation}\label{eq_bessel}
\lambda_nJ_1(\lambda_n)=\mathrm{Bi}\,J_0(\lambda_n),\qquad
\mathrm{Bi}=\frac{hR}{k}.
\end{equation}
由式\eqref{eq_bessel}确定各阶模态，结合$T_a(0)=28$，得到温度参考解
\begin{equation}\label{eq_duhamel}
T_{\mathrm{ref}}(r,t)=T_a(t)-\sum_{n=1}^{\infty}A_nJ_0(\lambda_nr/R)
\int_0^t e^{-\beta_n(t-\tau)}T_a'(\tau)\dd\tau,
\end{equation}
其中，$\beta_n=\alpha\lambda_n^2/R^2$，模态系数为
\[
A_n=\frac{2J_1(\lambda_n)}{\lambda_n[J_0^2(\lambda_n)+J_1^2(\lambda_n)]}.
\]
式\eqref{eq_duhamel}使用与有限体积计算相同的PCHIP环境输入，通过模态常微分方程计算卷积积分，不采用径向有限体积离散。完整版本附录C给出分离变量及Duhamel叠加推导。将级数项数由480增至960时，参考解的最大变化为$2.7632\times10^{-9}$ $^\circ$C；采用960项的参考解与有限体积解在全部输出点上的最大温度差为$2.1547\times10^{-7}$ $^\circ$C。

该项检验表明，在所采用的常系数一维径向导热模型及Robin边界下，温度空间离散和时间推进具有较高一致性。它不能验证忽略端面、忽略潜热和将$C_a$作为等效水分势等物理假设，因为这些假设同时会改变实际温度模型；相关影响仍需由模型适用性分析处理。

\medskip
\noindent\textbf{（3）非线性水分方程的辅助检验}

正式水分方程的$D$随$C$变化，不能直接使用常系数解析解作为其真值。为单独检查控制体几何、内部面通量和Robin边界的离散实现，另将$D$临时冻结为初始值$D(C_0)$，再与具有相同边界条件的Bessel展开比较。该辅助算例中，数值解与解析解的最大差为$1.5090\times10^{-6}$ kg/kg。由于冻结$D$删除了$D'(C)(\partial_rC)^2$项，该结果只能说明线性扩散部分的离散实现正确；正式非线性水分场仍由上述三级网格收敛检验和下述收支检验共同支持。

\medskip
\noindent\textbf{（4）热量与水分收支一致性}

对式\eqref{eq_fvm}在所有控制体上求和，内部面通量成对抵消，只保留表面交换项。定义有效蓄热量$\mathcal E_T(t)$和归一化水分存量$M_C(t)$，并分别以表面累计换热与累计排湿构造右侧量：
\begin{equation}\label{eq_q1_balance}
\begin{aligned}
\mathcal E_T(t)&=2\pi L\rho c_p\sum_iW_i[T_i(t)-T_i(0)],&
Q_T(t)&=2\pi RL\int_0^t h[T_a(\tau)-T_s(\tau)]\dd\tau,\\
M_C(t)&=\sum_iW_iC_i(t),&
L_C(t)&=R\int_0^t h_m[C_s(\tau)-C_a(\tau)]\dd\tau.
\end{aligned}
\end{equation}
式\eqref{eq_q1_balance}左右两侧在理论上应满足$\mathcal E_T(t)=Q_T(t)$及$M_C(t)-M_C(0)=-L_C(t)$。采用连续解的稠密输出在PCHIP分区间上作自适应积分，前1800 s的有效热收支相对残差为$2.0484\times10^{-9}$，归一化水分收支相对残差为$2.4254\times10^{-11}$。其中$M_C$仅用于检验含水率方程的归一化守恒；若要换算为实际水质量，还需乘以$2\pi L\rho_d$并明确参考干物质密度$\rho_d$。

上述收支检验主要说明内部面通量的符号、几何权重和表面边界处理与所建方程一致，并不等同于对真实药材失水量的实验验证。

\medskip
\noindent\textbf{（5）检验结论}

综合来看，温度场已经通过独立解析参考、空间和时间收敛以及热量收支三类检验；水分场受非线性扩散系数限制，没有可直接使用的解析真值，但通过了三级网格与时间收敛检验，其线性扩散部分另通过冻结系数算例和归一化收支检验。上述证据支持“本文数值解在所建立的一维等效方程下具有稳定的离散实现”，不能证明忽略端面、蒸发潜热和等效水分边界等物理简化已经达到实物验证精度。相关适用边界在后续模型评价中继续说明。
\FloatBarrier
